\documentclass[11pt,a4paper]{report}

\usepackage[utf8]{inputenc}
\usepackage[T1]{fontenc}
\usepackage[french]{babel}
\usepackage{lmodern}
\usepackage{microtype}
\usepackage[a4paper,margin=2.5cm,headheight=15pt]{geometry}
\usepackage{xcolor}
\usepackage{graphicx}
\usepackage{booktabs}
\usepackage{tabularx}
\usepackage{longtable}
\usepackage{array}
\usepackage{enumitem}
\usepackage{fancyhdr}
\usepackage{tikz}
\usetikzlibrary{arrows.meta,positioning,shapes.geometric,fit,calc}
\usepackage[hidelinks]{hyperref}

\definecolor{navy}{HTML}{17365D}
\definecolor{blue}{HTML}{2563EB}
\definecolor{lightblue}{HTML}{EAF2FF}
\definecolor{teal}{HTML}{0F766E}
\definecolor{lightteal}{HTML}{E6F5F2}
\definecolor{amber}{HTML}{A16207}
\definecolor{lightamber}{HTML}{FFF7E6}
\definecolor{gray}{HTML}{526174}
\definecolor{lightgray}{HTML}{F3F5F8}
\definecolor{darkred}{HTML}{9B2C2C}

\hypersetup{
  pdftitle={Rapport de stage -- Système de suivi des rendez-vous et du no-show},
  pdfauthor={Prénom NOM -- à compléter},
  pdfsubject={Rapport de stage en informatique},
  colorlinks=true,
  linkcolor=navy,
  urlcolor=blue,
  citecolor=teal
}

\setlength{\parskip}{0.45em}
\setlength{\parindent}{0pt}
\setlist[itemize]{leftmargin=1.5em,itemsep=0.2em,topsep=0.25em}
\setlist[enumerate]{leftmargin=1.7em,itemsep=0.25em,topsep=0.25em}
\renewcommand{\arraystretch}{1.18}
\newcommand{\code}[1]{\texttt{#1}}
\newcommand{\implemented}{\textcolor{teal}{\textbf{En place}}}
\newcommand{\planned}{\textcolor{amber}{\textbf{À réaliser / vérifier}}}
\newcommand{\placeholder}[1]{\textcolor{darkred}{\textbf{[#1]}}}

\pagestyle{fancy}
\fancyhf{}
\fancyhead[L]{\small Rapport de stage -- Système de suivi des rendez-vous}
\fancyhead[R]{\small Prototype logiciel}
\fancyfoot[C]{\thepage}
\renewcommand{\headrulewidth}{0.4pt}

\tikzset{
  component/.style={draw=navy,fill=lightblue,rounded corners=3pt,
    align=center,minimum height=1.0cm,text width=3.0cm,font=\small},
  storage/.style={draw=teal,fill=lightteal,rounded corners=3pt,
    align=center,minimum height=0.95cm,text width=2.7cm,font=\small},
  external/.style={draw=amber,fill=lightamber,rounded corners=3pt,
    align=center,minimum height=0.95cm,text width=2.7cm,font=\small},
  flow/.style={-{Stealth[length=2.2mm]},thick,draw=gray},
  flowlabel/.style={font=\scriptsize,fill=white,inner sep=1.5pt,align=center}
}

\begin{document}

\begin{titlepage}
  \centering
  \vspace*{1.4cm}
  {\Large\textsc{\placeholder{Nom de l'établissement}}\\[0.35em]
  \large\placeholder{Filière / département}\par}
  \vspace{2.0cm}
  {\color{blue}\rule{\textwidth}{1.2pt}\par}
  \vspace{0.6cm}
  {\Huge\bfseries Système de suivi des rendez-vous\\[0.25em]
  et d'estimation du risque de non-présentation\par}
  \vspace{0.55cm}
  {\Large Rapport de stage -- conception et réalisation d'un prototype\par}
  \vspace{0.6cm}
  {\color{blue}\rule{\textwidth}{1.2pt}\par}
  \vspace{1.8cm}
  \begin{tabular}{>{\bfseries}p{5.0cm}p{8.0cm}}
    Étudiant(e) & \placeholder{Prénom NOM}\\[0.65em]
    Établissement & \placeholder{Nom de l'école / université}\\[0.65em]
    Organisme d'accueil & \placeholder{Nom de l'organisme}\\[0.65em]
    Encadrant(e) en entreprise & \placeholder{Nom et fonction}\\[0.65em]
    Encadrant(e) pédagogique & \placeholder{Nom et fonction}\\[0.65em]
    Période du stage & \placeholder{Date de début -- date de fin}\\[0.65em]
    Année universitaire & \placeholder{AAAA--AAAA}
  \end{tabular}
  \vfill
  {\large \placeholder{Ville}, \placeholder{mois et année de remise}\par}
  \vspace{0.7cm}
\end{titlepage}

\pagenumbering{roman}
\chapter*{Résumé}
\addcontentsline{toc}{chapter}{Résumé}
Ce rapport décrit un prototype logiciel de gestion et d'analyse de rendez-vous,
centré sur l'estimation expérimentale du risque de non-présentation
(\emph{no-show}). Le dépôt étudié assemble une API REST écrite avec FastAPI, une
base relationnelle PostgreSQL accessible par SQLAlchemy, une interface
analytique Streamlit, un classifieur XGBoost chargé depuis un artefact Joblib et
un assistant de suggestions utilisant un service Ollama local.

Le code comporte notamment des opérations de consultation et de mise à jour des
patients, médecins et rendez-vous, une authentification par mot de passe haché
et jeton JWT, un calcul de trois variables temporelles et un tableau de bord
permettant de consulter des indicateurs. Il comprend également des tests
unitaires et des contrôles de validation. Ces éléments constituent un état
technique de prototype, et non une solution clinique validée ou une preuve de
conformité réglementaire. Les données générées par le dépôt sont synthétiques ;
aucun résultat de performance réel n'est fourni par le code examiné.

\textbf{Mots-clés :} rendez-vous, no-show, FastAPI, Streamlit, PostgreSQL,
apprentissage automatique, protection des données.

\chapter*{Note de lecture et hypothèses}
\addcontentsline{toc}{chapter}{Note de lecture et hypothèses}
Le rapport est établi à partir des fichiers présents dans le dépôt
\code{healthcare-noshow-system} au moment de son examen. Il ne présume ni du
contexte réel du stage ni d'un déploiement en établissement de santé. Les champs
rouges de la page de garde sont à renseigner par l'étudiant. Les résultats
mesurés, l'environnement d'exécution réellement utilisé, les choix
d'architecture validés par les encadrants et les éventuelles évolutions
postérieures au dépôt doivent être complétés à partir de preuves vérifiables.

Dans la suite, les éléments qualifiés d'« en place » sont repérés dans le code.
Les éléments qualifiés d'« à réaliser / vérifier » sont des recommandations,
des objectifs de feuille de route ou des obligations d'évaluation qui ne doivent
pas être confondus avec une fonction déjà livrée. Aucun dossier patient réel
n'a été consulté pour élaborer ce rapport.

\tableofcontents
\clearpage
\pagenumbering{arabic}

\chapter{Introduction et contexte}
\section{Problématique}
Les rendez-vous manqués peuvent compliquer l'organisation d'un service de
santé. Un outil de suivi peut aider le personnel à repérer les rendez-vous à
examiner et à proposer des actions organisationnelles telles qu'un rappel ou
une demande de confirmation. Une telle aide ne remplace ni l'évaluation par le
personnel ni une décision clinique.

Le projet étudié réunit une gestion élémentaire de données de rendez-vous et
une estimation de risque expérimentale. Son objectif technique observable est
de relier une API de persistance et de prédiction à une interface de
visualisation, puis d'exposer, sur demande, des suggestions générées par un
modèle de langage exécuté localement.

\section{Objectifs du prototype}
\begin{itemize}
  \item gérer les enregistrements des patients, des médecins et des rendez-vous ;
  \item contrôler les accès à l'API au moyen d'une authentification par jeton ;
  \item calculer et exposer une estimation de probabilité de no-show, seule ou
        par lot ;
  \item visualiser les rendez-vous d'une date, des statistiques descriptives
        et un historique patient dans Streamlit ;
  \item expérimenter des suggestions de suivi non médicales avec Ollama local ;
  \item expliciter les limites de sécurité, de confidentialité et de validation
        avant toute perspective de mise en service.
\end{itemize}

\section{Périmètre et état des preuves}
\begin{tabularx}{\textwidth}{>{\bfseries}p{3.1cm}X}
\toprule
Sujet & État constaté dans le dépôt\\
\midrule
API et données & \implemented{} CRUD partiel des entités principales,
validation Pydantic et couche SQLAlchemy.\\
Authentification & \implemented{} connexion par formulaire OAuth2, mot de
passe bcrypt et JWT signé à expiration.\\
Interface & \implemented{} tableau de bord Streamlit connecté à l'API,
consultation et graphiques.\\
Modèle & \implemented{} artefact pré-entraîné et script de réentraînement ;
performance clinique \planned{}.\\
Assistant & \implemented{} appel Ollama limité à une adresse locale et à des
variables numériques ; validation opérationnelle et conformité \planned{}.\\
Réglementation & Analyse préparatoire seulement ; aucune conformité CNDP ou
à la loi marocaine 09-08 n'est revendiquée.\\
\bottomrule
\end{tabularx}

\chapter{Architecture et composants}
\section{Vue d'ensemble}
Le dépôt est organisé en modules Python plutôt qu'en paquet applicatif
complexe. \code{main.py} crée l'application FastAPI, charge
\code{model.pkl}, déclare les routes et assemble leurs dépendances.
\code{models.py} décrit les tables SQLAlchemy, \code{schemas.py} les contrats
Pydantic et \code{database.py} la configuration de la connexion. Le tableau de
bord (\code{dashboard.py}) consomme l'API via HTTP. Le script
\code{train_model.py} et le module \code{feature_engineering.py} portent le
pipeline de données et de modèle. \code{assistant_agent.py} encapsule la
requête au service Ollama.

\begin{figure}[htbp]
\centering
\resizebox{\textwidth}{!}{%
\begin{tikzpicture}[node distance=1.1cm and 1.15cm]
  \node[component] (staff) {Personnel\\(navigateur)};
  \node[component,right=of staff] (dash) {Interface Streamlit\\
    \code{dashboard.py}};
  \node[component,right=of dash] (api) {API REST FastAPI\\
    \code{main.py}};
  \node[storage,right=of api] (db) {PostgreSQL\\patients, médecins,\\
    rendez-vous, comptes, audit};
  \node[component,below=1.1cm of api] (model) {Inférence\\
    \code{model.pkl} / XGBoost};
  \node[external,below=1.1cm of db] (ollama) {Ollama local\\
    modèle de langage};
  \node[component,below=1.1cm of dash] (training) {Préparation et entraînement\\
    \code{feature\_engineering.py}\\\code{train\_model.py}};

  \draw[flow] (staff) -- (dash);
  \draw[flow] (dash) -- node[flowlabel,above]{HTTP / JSON\\Bearer JWT} (api);
  \draw[flow] (api) -- node[flowlabel,above]{SQLAlchemy} (db);
  \draw[flow] (api) -- node[flowlabel,right]{features\\numériques} (model);
  \draw[flow] (api) -- node[flowlabel,right]{requête locale\\numérique} (ollama);
  \draw[flow] (training) -- node[flowlabel,above]{lecture SQL\\et sauvegarde} (db);
  \draw[flow] (training) -- node[flowlabel,left]{artefact modèle\\à déployer} (model);
\end{tikzpicture}}
\caption{Architecture logique observée dans le dépôt. La base et le service
Ollama doivent être configurés séparément ; le schéma ne représente pas un
déploiement de production validé.}
\label{fig:architecture}
\end{figure}

\section{Échanges et frontières de confiance}
Le tableau de bord ne dialogue pas directement avec PostgreSQL : il envoie des
requêtes à l'API et détient le jeton pendant la session Streamlit. L'API lit ou
écrit les dossiers via SQLAlchemy, puis réalise l'inférence sur des variables
numériques. Lorsqu'une suggestion est demandée, l'API relaie au processus
Ollama uniquement le score, sa classe de risque, le délai de réservation, le
jour de la semaine et l'historique no-show. Le nom, le téléphone, l'adresse
électronique et le motif de visite ne sont pas inclus dans le message adressé
au modèle de langage.

\begin{figure}[htbp]
\centering
\resizebox{0.96\textwidth}{!}{%
\begin{tikzpicture}[node distance=0.72cm and 0.8cm]
  \node[component,text width=2.7cm] (user) {Personnel\\choisit une date};
  \node[component,right=of user,text width=3.0cm] (dash) {Streamlit\\
    charge les pages d'API};
  \node[storage,right=of dash,text width=2.8cm] (db) {API + PostgreSQL\\
    dossiers et rendez-vous};
  \node[component,below=of db,text width=2.8cm] (features) {Variables\\
    délai, jour, antécédents};
  \node[component,left=of features,text width=2.8cm] (risk) {API + modèle\\
    score et classe};
  \node[component,left=of risk,text width=2.8cm] (display) {Tableau de bord\\
    liste et indicateurs};
  \node[external,below=of risk,text width=2.8cm] (genai) {Ollama local\\
    suggestion structurée};
  \node[component,below=of display,text width=2.8cm] (human) {Personnel\\
    valide ou ignore};

  \draw[flow] (user) -- (dash);
  \draw[flow] (dash) -- node[flowlabel,above]{Bearer + GET} (db);
  \draw[flow] (db) -- (features);
  \draw[flow] (features) -- (risk);
  \draw[flow] (risk) -- (display);
  \draw[flow] (display) -- (human);
  \draw[flow] (risk) -- node[flowlabel,right]{sur action explicite} (genai);
  \draw[flow] (genai) -- (human);
\end{tikzpicture}}
\caption{Parcours d'une consultation et d'une suggestion. Le diagramme
représente la logique codée, pas une décision automatisée : toute suggestion
est destinée à une revue humaine.}
\label{fig:flux}
\end{figure}

\section{Technologies et lancement}
Les dépendances épinglées dans \code{requirements.txt} incluent FastAPI,
Uvicorn, SQLAlchemy, Alembic, psycopg2, Pydantic, Streamlit, Plotly, pandas,
scikit-learn, XGBoost, Joblib, bcrypt, PyJWT et requests. Le dépôt demande
Python 3.10 ou plus récent et PostgreSQL. Les variables de configuration
figurent comme exemples dans \code{.env.example}; le fichier \code{.env} ne
doit pas être publié.

Le README indique les commandes de démarrage suivantes après installation,
configuration et application de la migration :
\begin{itemize}
  \item \code{alembic upgrade head} pour créer les tables absentes ;
  \item \code{python create_user.py} pour créer un compte initial ;
  \item \code{uvicorn main:app --reload} pour démarrer l'API ;
  \item \code{streamlit run dashboard.py} pour démarrer le tableau de bord.
\end{itemize}
L'URL par défaut de l'API est \url{http://127.0.0.1:8000}; la documentation
OpenAPI est servie par FastAPI à \url{/docs}. Le mode \code{--reload} est un
mode de développement, non une recommandation de déploiement.

\chapter{API, authentification et persistance}
\section{API FastAPI}
Les requêtes sont validées avec les schémas Pydantic. Les routes sauf
\code{POST /auth/login} sont enregistrées sous un routeur protégé par
\code{get\_current\_user}. Les listes acceptent une pagination
\code{limit}/\code{offset}; les listes patients et médecins ont une limite par
défaut de 100 et maximale de 500. La liste des rendez-vous accepte aussi
\code{from\_date} et \code{to\_date}. Une plage inversée est rejetée.

\begin{longtable}{>{\raggedright\arraybackslash}p{2.35cm}
                  >{\raggedright\arraybackslash}p{4.3cm}
                  >{\raggedright\arraybackslash}p{7.1cm}}
\toprule
Méthode & Route & Fonction et restrictions constatées\\
\midrule
\endhead
POST & \code{/auth/login} & Connexion par formulaire OAuth2 ; retourne un
jeton bearer ou une réponse 401. Route publique nécessaire à l'obtention du
jeton.\\
POST & \code{/predict} & Estimation pour un rendez-vous à partir de
\code{delai\_jours}, \code{jour\_semaine} et
\code{nb\_no\_show\_precedents}.\\
POST & \code{/predict/batch} & Inférence par lot ; entre 1 et 500 éléments,
un résultat par élément.\\
POST & \code{/assistant/no-show} & Génère une suggestion de suivi à partir du
score, de sa classe et des trois variables numériques ; accès protégé.\\
POST / GET / GET par ID / PUT & \code{/patients} et
\code{/patients/\{patient\_id\}} & Création, liste paginée, lecture et mise à
jour d'un patient. Aucune route DELETE n'est définie.\\
POST / GET / GET par ID / PUT & \code{/doctors} et
\code{/doctors/\{doctor\_id\}} & Création, liste paginée, lecture et mise à
jour d'un médecin. Aucune route DELETE n'est définie.\\
POST / GET / GET par ID / PUT & \code{/appointments} et
\code{/appointments/\{appointment\_id\}} & Création, liste paginée et filtrée
par date, lecture et mise à jour d'un rendez-vous. Aucune route DELETE n'est
définie.\\
\bottomrule
\end{longtable}

Les schémas imposent notamment des formats d'adresse électronique, des
longueurs de champs, une date de naissance non future, une date de rendez-vous
non antérieure à la réservation et un statut appartenant à une liste codée
(\code{no\_show}, \code{confirmé}, \code{confirmed}, \code{scheduled},
\code{cancelled}, \code{annulé}, \code{présent}, \code{attended},
\code{completed}). Les paramètres de prédiction sont bornés et le niveau de
risque fourni à l'assistant doit correspondre au score suivant les seuils
définis dans \code{risk.py}.

\section{Authentification et audit}
Le mot de passe est stocké sous forme de hachage bcrypt. Le script
\code{create\_user.py} valide un nom d'utilisateur, impose au moins huit
caractères au mot de passe (et au plus 72 octets UTF-8) et accepte les rôles
\code{admin} ou \code{staff}. L'API retourne un JWT signé en HS256, portant le
sujet utilisateur et une date d'expiration (60 minutes par défaut, configurable
par \code{ACCESS\_TOKEN\_EXPIRE\_MINUTES}). Une clé \code{SECRET\_KEY} d'au
moins 32 caractères est requise au démarrage.

La dépendance d'authentification vérifie le jeton, retrouve l'utilisateur en
base et écrit un enregistrement d'audit avec acteur, méthode HTTP, gabarit de
route et horodatage. Le code n'enregistre pas le contenu de la requête dans
cet audit. L'authentification n'implémente cependant pas de permissions
distinctes par rôle : les rôles sont reconnus comme valides, mais
\code{admin} et \code{staff} accèdent au même ensemble de routes. Les échecs
de connexion ne sont pas audités ni limités en débit dans le code examiné.

\section{Schéma relationnel}
Le modèle SQLAlchemy contient cinq tables applicatives. Les types et
contraintes ci-dessous sont ceux des modèles et de la migration initiale.

\begin{figure}[htbp]
\centering
\resizebox{0.92\textwidth}{!}{%
\begin{tikzpicture}[node distance=1.5cm and 1.5cm]
  \node[storage,text width=4.2cm,minimum height=1.45cm] (patients)
    {\textbf{patients}\\PK id\\name, telephone, email, date\_naissance};
  \node[storage,right=of patients,text width=4.2cm,minimum height=1.45cm] (appointments)
    {\textbf{appointments}\\PK id\\FK patient\_id, FK doctor\_id\\dates, motif\_visite, statut};
  \node[storage,right=of appointments,text width=4.0cm,minimum height=1.45cm] (doctors)
    {\textbf{doctors}\\PK id\\name, telephone, email, departement};
  \node[external,below=1.4cm of appointments,text width=4.0cm,minimum height=1.35cm] (audit)
    {\textbf{audit\_logs}\\PK id, FK actor\_id nullable\\action, entity, created\_at};
  \node[external,below=1.4cm of patients,text width=4.0cm,minimum height=1.35cm] (users)
    {\textbf{users}\\PK id, username unique\\hashed\_password, role};

  \draw[flow] (patients.east) -- node[flowlabel,above]{1 vers 0..n} (appointments.west);
  \draw[flow] (doctors.west) -- node[flowlabel,above]{1 vers 0..n} (appointments.east);
  \draw[flow] (users.east) -- node[flowlabel,above,sloped]{1 vers 0..n} (audit.south west);
\end{tikzpicture}}
\caption{Relations fonctionnelles définies par les clés étrangères. Les comptes
et journaux ne sont pas reliés aux rendez-vous.}
\label{fig:schema}
\end{figure}

\begin{longtable}{>{\raggedright\arraybackslash}p{2.6cm}
                  >{\raggedright\arraybackslash}p{4.3cm}
                  >{\raggedright\arraybackslash}p{6.9cm}}
\toprule
Table & Attributs principaux & Rôle et remarques\\
\midrule
\endhead
\code{patients} & \code{id} PK, \code{name}, \code{telephone},
\code{email}, \code{date\_naissance} & Informations d'identification et de
contact ; les champs sont obligatoires.\\
\code{doctors} & \code{id} PK, \code{name}, \code{telephone},
\code{email}, \code{departement} & Informations d'identification et
d'affectation du médecin.\\
\code{appointments} & \code{id} PK, \code{patient\_id} FK,
\code{doctor\_id} FK, \code{date\_reservation}, \code{date\_rdv},
\code{motif\_visite}, \code{statut} & Rendez-vous relié à un patient et un
médecin. Index sur les deux clés étrangères et la date du rendez-vous.\\
\code{users} & \code{id} PK, \code{username} unique,
\code{hashed\_password}, \code{role} & Comptes applicatifs ; contrainte SQL
sur les rôles \code{admin} et \code{staff}.\\
\code{audit\_logs} & \code{id} PK, \code{actor\_id} FK nullable,
\code{action}, \code{entity}, \code{created\_at} & Trace les accès
authentifiés ; index sur l'acteur et l'horodatage ; suppression du compte
associé préserve le journal en mettant l'acteur à NULL.\\
\bottomrule
\end{longtable}

La configuration accepte soit une \code{DATABASE\_URL}, soit les variables
\code{DB\_USER}, \code{DB\_PASSWORD}, \code{DB\_HOST}, \code{DB\_PORT} et
\code{DB\_NAME}. Le constructeur SQLAlchemy compose les paramètres séparément
pour éviter les problèmes d'échappement du mot de passe. Alembic applique une
migration initiale qui crée les tables manquantes et vérifie les index. Elle
ne modifie pas les colonnes d'une table déjà présente et est volontairement
irréversible afin de protéger les données existantes ; les divergences de
schéma héritées exigent donc une migration explicite.

\chapter{Tableau de bord Streamlit}
\section{Fonctions visibles}
L'interface fournit un formulaire de connexion, une déconnexion, un filtre
par date, une liste des rendez-vous de la journée enrichie par des scores,
trois indicateurs (volume de rendez-vous, nombre de risques élevés, risque
moyen), ainsi que deux vues statistiques : le taux mensuel de no-show et le
taux par motif. Un onglet de recherche patient présente le nombre de no-shows
et l'historique des rendez-vous de ce patient. Une action explicite ouvre
l'assistant pour un rendez-vous sélectionné.

Les requêtes HTTP utilisent le jeton bearer conservé dans l'état de session.
Le client a un délai d'attente de dix secondes par défaut et traite une
réponse 401 en effaçant le jeton. L'appel GenAI utilise un délai de 90 secondes.
L'interface charge les rendez-vous et les patients par pages de 500, puis
construit des DataFrames en mémoire. Pour la date sélectionnée, elle envoie
les variables d'inférence en lots à \code{/predict/batch}.

\section{Limites d'usage et d'échelle}
Cette interface est un outil de consultation, pas une application complète
d'administration : aucun formulaire de saisie ou de suppression n'est codé
dans Streamlit, et les opérations de création ou de mise à jour sont
uniquement disponibles comme routes d'API. Malgré la pagination HTTP, le
tableau de bord finit par récupérer l'ensemble des pages et de l'historique ;
les statistiques sont calculées côté client. Le filtrage principal des
rendez-vous par date se fait après téléchargement. Pour un grand volume, il
faudrait déplacer les agrégations et filtres vers la base et limiter
l'historique consulté.

Le tableau de bord manipule en mémoire des identifiants, noms et motifs de
visite pour afficher ses vues, même si ces valeurs ne sont pas transmises à
l'assistant GenAI. Il ne doit donc pas être traité comme une interface déjà
autorisée à recevoir de vraies données de santé.

\chapter{Données et pipeline d'apprentissage automatique}
\section{Variables et construction des caractéristiques}
Le module \code{feature\_engineering.py} valide la présence des identifiants,
dates et statuts, rejette les valeurs manquantes et les dates de rendez-vous
antérieures à la réservation, puis trie les lignes par patient, date et
identifiant. Les trois caractéristiques utilisées par le modèle sont :
\begin{itemize}
  \item \code{delai\_jours} : nombre de jours entre réservation et rendez-vous ;
  \item \code{jour\_semaine} : jour de la semaine du rendez-vous (0 à 6) ;
  \item \code{nb\_no\_show\_precedents} : nombre de no-shows du patient aux
        dates strictement antérieures.
\end{itemize}
Pour éviter qu'un rendez-vous n'influence les autres rendez-vous du même jour,
les no-shows sont agrégés par patient et par date, puis le cumul est décalé
d'une date avant d'être associé aux lignes du jour. La cible est binaire :
statut \code{no\_show} contre les statuts d'assiduité observée
(\code{confirmé}, \code{confirmed}, \code{présent}, \code{attended},
\code{completed}). Les statuts d'annulation et les rendez-vous encore
programmés ne servent pas de négatifs au réentraînement.

\section{Entraînement et inférence}
\code{train\_model.py} lit dans PostgreSQL les colonnes d'identifiant patient,
dates et statut, construit les variables, puis entraîne un
\code{XGBClassifier} à 200 estimateurs avec une séparation chronologique
approximativement 80/20 sur les dates distinctes. Le code vérifie que les
partitions sont exploitables et que l'entraînement contient les deux classes.
Il calcule accuracy, précision, rappel, F1, ROC-AUC lorsque cette métrique est
définissable, PR-AUC et score de Brier, puis écrit le modèle dans
\code{model.pkl}. L'inférence dans \code{main.py} retourne une classe prédite,
une probabilité arrondie à quatre décimales et une bande de risque :
faible si score inférieur à 0,3 ; moyen de 0,3 inclus à 0,6 exclu ; élevé à
partir de 0,6.

Le fichier \code{model.pkl} présent dans le dépôt est un artefact pré-entraîné.
Le code ne fournit pas, à lui seul, de provenance complète de cet artefact,
de versionnement des jeux de données ou de rapport de validation clinique.
Le script de génération produit des rendez-vous synthétiques de façon
déterministe (graine par défaut 42) ; il insère sans supprimer les données
existantes sauf option explicite \code{--replace}. Cette option tronque les
rendez-vous et est destructive ; elle n'a pas été exécutée dans le cadre de
ce rapport.

\section{Résultats : champs à compléter, sans extrapolation}
Aucun score mesuré n'est affirmé ici. Les métriques imprimées par
\code{train\_model.py} doivent être calculées sur un jeu de données et un
protocole documentés, puis reportées avec leurs limites. L'évaluation sur
données synthétiques ne démontrerait pas les performances sur une population
réelle.

\begin{tabularx}{\textwidth}{>{\bfseries}p{3.1cm}X}
\toprule
Élément du compte rendu & Valeur vérifiée à compléter par l'étudiant\\
\midrule
Origine, période et effectif du jeu évalué &
\placeholder{À documenter ; préciser données synthétiques ou autorisées}\\
Date de séparation et nombres train/test &
\placeholder{Valeurs issues d'une exécution reproductible}\\
Accuracy, précision, rappel et F1 &
\placeholder{Résultats mesurés avec définition du seuil}\\
ROC-AUC, PR-AUC et score de Brier &
\placeholder{Résultats mesurés ; préciser incertitudes et déséquilibre}\\
Évaluation par sous-groupes et impact opérationnel &
\placeholder{Non disponible dans le dépôt ; à mesurer avant tout usage}\\
\bottomrule
\end{tabularx}

\chapter{Assistant génératif local}
\section{Implémentation observée}
Le module \code{assistant\_agent.py} construit un appel au point
\code{/api/chat} d'Ollama. Par défaut, le modèle est \code{qwen2.5:3b} et le
service est configuré sur \code{http://127.0.0.1:11434}. Le code refuse une
URL qui n'utilise pas HTTP sur une adresse de boucle locale, ou qui contient
des informations d'identification, une query ou un fragment. Il désactive les
redirections et fixe un délai maximal de 90 secondes.

Le prompt demande une réponse en français au format JSON avec un résumé et
une à trois actions choisies parmi quatre codes : rappel neutre, demande de
confirmation, appel du personnel ou absence d'action supplémentaire. La
réponse est validée par Pydantic ; les codes autorisés sont transformés en
libellés compréhensibles et accompagnés d'un avertissement. En cas d'indisponibilité
du service, de réponse HTTP non valide ou de contenu mal formé, l'API renvoie
une erreur de service appropriée.

\section{Limites}
« Local » décrit ici la destination configurée de l'appel, non une garantie
de sécurité, de confidentialité ou de conformité : le contrôle de boucle
locale ne couvre pas à lui seul la machine hôte, l'accès local au service, le
téléchargement du modèle, la journalisation, la rétention des prompts ou les
composants système. L'interface et l'API doivent rester sous contrôle d'un
opérateur, et toute suggestion doit faire l'objet d'une validation humaine.
L'assistant est conçu pour des suggestions d'organisation, pas pour un
diagnostic, une décision de soin ou une mesure punitive. Les tests simulent
l'appel HTTP ; ils ne constituent pas une validation métier des productions
d'un modèle Ollama réel.

\chapter{Sécurité, confidentialité et limites réglementaires}
\section{Contrôles présents}
Le dépôt contient des mesures utiles au stade du prototype : hachage bcrypt,
jetons signés et expirants, validation des données d'entrée, pagination
bornée, journalisation des requêtes authentifiées sans corps de requête,
restrictions locales pour Ollama et exclusion de \code{.env} de Git. Le
générateur de données ne supprime pas les rendez-vous par défaut.

\section{Contrôles manquants ou à confirmer}
Les éléments suivants ne sont pas attestés comme configurés ou éprouvés par
le seul code :
\begin{itemize}
  \item HTTPS de bout en bout, chiffrement et contrôle d'accès de la base,
        sauvegardes chiffrées et restauration testée ;
  \item limitation des tentatives de connexion, MFA, révocation de session,
        cycle de vie des comptes et permissions différentes selon les rôles ;
  \item politique de conservation et suppression des données et des journaux,
        supervision, alertes, gestion d'incident et audit d'intégrité ;
  \item information des personnes, gestion de leurs droits, minimisation
        validée et procédures d'archivage ou de correction ;
  \item analyse formelle des risques, inventaire des sous-traitants, de
        l'hébergement, des accès de support et d'éventuels transferts ;
  \item tests de sécurité et validation de configuration pour un environnement
        de production.
\end{itemize}

\section{CNDP et loi marocaine 09-08 : prudence nécessaire}
Le dépôt comporte un document préparatoire \code{SECURITY\_PRIVACY.md} qui
rappelle que ce prototype ne constitue pas une déclaration de conformité.
Les informations liées à la santé nécessitent une analyse juridique et de
gouvernance spécifique. Le responsable de traitement, les finalités, la base
juridique, les formalités applicables auprès de la CNDP, l'information des
personnes, leurs droits, la durée de conservation et les transferts éventuels
doivent être déterminés par les personnes compétentes et confirmés à partir
des règles en vigueur. Ce rapport n'apporte pas de conseil juridique et ne
conclut pas que le système satisfait à la loi 09-08 ou à une exigence de la
CNDP.

Avant toute utilisation de données réelles, il faut obtenir les validations
institutionnelles et juridiques requises, réaliser une analyse de risques,
mettre en œuvre des protections techniques et organisationnelles adaptées
et documenter les responsabilités. L'exécution locale d'un modèle ne
dispense pas de ces vérifications. Aucune donnée patient réelle n'a été
consultée pour préparer ce rapport.

\section{Performances et équité : pas de validation clinique}
Les scores sont produits par un modèle expérimental et ne doivent pas être
utilisés seuls pour refuser, retarder ou prioriser l'accès aux soins. Le dépôt
ne fournit aucune preuve d'évaluation prospective, de calibration sur données
représentatives, de mesure des disparités entre sous-groupes, d'évaluation
d'impact ou de validation clinique. Les classes de risque sont des seuils
techniques configurés, non des niveaux approuvés par une autorité médicale.
La vérification des étiquettes, des rendez-vous dupliqués, des annulations et
des données manquantes reste nécessaire avant de tirer des conclusions.

\chapter{Tests, qualité et reproductibilité}
\section{Tests présents}
Le fichier \code{tests/test\_project.py} couvre notamment :
\begin{itemize}
  \item l'ingénierie des variables, l'exclusion des événements du même jour,
        l'ordre des colonnes et le rejet de données incohérentes ;
  \item les bornes des entrées Pydantic, les statuts et dates de rendez-vous,
        les seuils de classe de risque et la validation d'adresse électronique ;
  \item le découpage chronologique et le calcul des métriques d'entraînement,
        ainsi que l'exclusion des rendez-vous annulés ;
  \item les données envoyées à Ollama, le refus d'une URL distante et la
        validation de la réponse générée ;
  \item la création non destructive des tables manquantes par migration et
        des cas d'API : authentification, prédiction, lots et audit.
\end{itemize}
Les tests utilisent une base SQLite en mémoire ou de test, des données
fabriquées et des appels externes simulés. Le workflow GitHub Actions
\code{.github/workflows/tests.yml} installe les dépendances, compile les
modules Python, exécute la suite unittest et vérifie la découverte de
l'historique Alembic sur les événements push et pull request.

\section{Ce que les tests ne prouvent pas}
Les tests unitaires ne démontrent pas le fonctionnement sur PostgreSQL
configuré pour un déploiement, la robustesse sous charge, la sécurité d'une
configuration réseau, la qualité des prédictions sur des données de terrain
ou la conformité réglementaire. Ils ne remplacent pas des tests d'intégration
et de bout en bout, des tests d'autorisation par rôle, une analyse de
dépendances ou des essais utilisateur. L'état des tests à la date de remise
doit être renseigné à partir d'une exécution reproductible :
\[
\text{Commande : } \code{python -m unittest discover -s tests -v}
\qquad
\text{Résultat observé : } \placeholder{à renseigner}
\]

\chapter{Feuille de route et recommandations}
Le fichier \code{IMPROVEMENT\_ROADMAP.md} organise les évolutions en priorités.
Le tableau suivant distingue les recommandations de l'état implémenté.

\begin{longtable}{>{\raggedright\arraybackslash}p{1.5cm}
                  >{\raggedright\arraybackslash}p{5.5cm}
                  >{\raggedright\arraybackslash}p{6.8cm}}
\toprule
Priorité & Travaux recommandés & Résultat attendu\\
\midrule
\endhead
P0 & \planned{} Revoir l'autorisation de bout en bout, les rôles, la loi
09-08/CNDP, le responsable et les finalités ; exiger HTTPS, gestion sûre des
secrets, chiffrement des connexions et sauvegardes ; ajouter limitation des
tentatives, MFA, gestion des comptes, rétention d'audit et restauration
testée. & Réduire les risques de lancement et établir les conditions de
traitement légitime et sécurisé avant toute donnée réelle.\\
P1 & \planned{} Déplacer filtres et agrégations vers la base, borner
l'historique, vérifier les plans de requête et index, ajouter santé et
disponibilité, erreurs structurées, suivi, traçage, CI élargie et
environnements reproductibles. & Améliorer la fiabilité, la performance et
l'exploitabilité du service.\\
P2 & \planned{} Versionner les jeux, variables, modèles et évaluations ;
documenter les étiquettes ; évaluer chronologiquement et par groupes,
calibration, impact opérationnel et dérive ; prévoir une revue humaine et
une correction des données. & Établir si le modèle apporte une aide mesurable,
équitable et acceptable, sans automatiser une décision de soin.\\
P3 & \planned{} Ajouter suppression/archivage avec droits et audit,
recherche et filtres plus complets, accessibilité et localisation, versionnage
API, exemples OpenAPI, journalisation prudente, tests d'intégration et de
charge sur données synthétiques. & Faire évoluer le prototype vers une
application maintenable et utilisable.\\
\bottomrule
\end{longtable}

\chapter{Conclusion}
Le dépôt met en œuvre une chaîne logicielle cohérente pour expérimenter une
application de rendez-vous : API FastAPI, stockage relationnel, authentification
JWT, tableau de bord Streamlit, calcul de caractéristiques temporelles,
inférence XGBoost et suggestions de workflow via Ollama local. Les migrations,
les validations de schéma et des tests unitaires donnent une base de
développement identifiable.

Cette cohérence technique ne suffit pas à autoriser un usage opérationnel
dans un environnement de soins. Les performances n'ont pas été établies sur
des données réelles représentatives, le système ne comporte pas toutes les
protections ni procédures nécessaires, et aucune conformité à la CNDP ou à la
loi 09-08 n'est déclarée. La prochaine étape est donc une phase de cadrage,
d'évaluation et de conformité, avant tout traitement de données personnelles
réelles ou toute utilisation de scores pour orienter une action.

\chapter*{Références techniques}
\addcontentsline{toc}{chapter}{Références techniques}
Les références principales utilisées pour décrire l'implémentation sont les
fichiers du dépôt :
\begin{itemize}
  \item \code{README.md} : installation, lancement, utilisation du modèle et
        de l'assistant ;
  \item \code{main.py}, \code{auth.py}, \code{schemas.py} : routes, sécurité
        applicative et validation ;
  \item \code{models.py}, \code{database.py}, \code{migrations/versions/} :
        schéma relationnel, connexion et migration ;
  \item \code{dashboard.py} : parcours de consultation et visualisations ;
  \item \code{feature\_engineering.py}, \code{train\_model.py},
        \code{risk.py}, \code{model.pkl} : variables, entraînement, seuils
        et artefact d'inférence ;
  \item \code{assistant\_agent.py} : intégration Ollama ;
  \item \code{tests/test\_project.py} et
        \code{.github/workflows/tests.yml} : tests et automatisation ;
  \item \code{SECURITY\_PRIVACY.md} et
        \code{IMPROVEMENT\_ROADMAP.md} : limites et travaux à prévoir.
\end{itemize}
Pour les obligations applicables, consulter les ressources officielles de la
CNDP (\url{https://www.cndp.ma/}) et solliciter les compétences juridiques et
institutionnelles pertinentes. La liste ci-dessus n'est pas une validation
juridique.

\appendix
\chapter{Checklist de personnalisation avant remise}
\begin{enumerate}
  \item Remplacer tous les champs \placeholder{entre crochets} de la page de
        garde par des informations exactes ; ne pas inventer le nom de
        l'établissement, les personnes ou les dates.
  \item Ajouter, si disponible, le contexte du stage réellement autorisé et
        les décisions confirmées par les encadrants, sans exposer de données
        personnelles ou de secret.
  \item Compléter les mesures uniquement après exécution et archiver le
        protocole, la source des données et le résultat reproductible.
  \item Maintenir la distinction entre fonctions présentes, évolutions
        planifiées et validations réglementaires encore requises.
  \item Relire le rapport avec l'encadrant et vérifier la politique de
        confidentialité et de diffusion de l'établissement d'accueil.
\end{enumerate}

\end{document}
